% =============================================================================
%  Tipos de vocabulário e computação em tempo de compilação
%  Cobre os subprojetos vocabulary_types, compile_time, cpp23 e modules.
% =============================================================================

% =============================================================================
\section{Tipos de vocabulário e computação em tempo de compilação}
\label{sec:vocab-compile-time}
% =============================================================================

Este capítulo trata de duas ideias que parecem não ter relação: um punhado de
tipos da biblioteca padrão --- \texttt{optional}, \texttt{variant},
\texttt{any}, \texttt{tuple}, \texttt{span} --- e um conjunto de palavras-chave
que mandam o compilador calcular coisas antes de o programa existir ---
\texttt{constexpr}, \texttt{consteval}, \texttt{constinit},
\texttt{static\_assert}.

A relação é esta: as duas movem informação para dentro do \textbf{tipo}, e tudo
o que está no tipo é verificado por alguém que nunca se distrai. Um
\texttt{optional<int>} diz, na assinatura, que pode não haver valor --- e quem
esquecer de testar recebe um erro em vez de ler lixo. Uma
\texttt{Matriz<2,3>} diz, no tipo, quantas linhas e colunas tem --- e quem
multiplicar por uma \texttt{Matriz<2,2>} não compila. É a mesma estratégia em
escalas diferentes: \textbf{fazer o compilador trabalhar no seu lugar}.

O capítulo percorre quatro subprojetos. \projeto{vocabulary\_types} apresenta os
tipos que atravessam fronteiras de biblioteca (\cref{sec:vocab-fundo});
\projeto{compile\_time} mostra o que o compilador calcula sozinho
(\cref{sec:ct-fundo}); \projeto{cpp23} \emph{mede} o que o compilador de
referência realmente implementa, em vez de supor (\cref{sec:cpp23-medido}); e
\projeto{modules} enfrenta a substituição do \texttt{\#include} e documenta dois
defeitos do GCC 13 encontrados na prática (\cref{sec:modulos}). Os dois
primeiros assuntos apareceram resumidos no \cref{sec:vocabulary} e no
\cref{sec:compile-time}; aqui são desenvolvidos do zero.

% =============================================================================
\subsection{O vocabulário compartilhado}
\label{sec:vocab-fundo}
% =============================================================================

Dois termos, antes de tudo. Uma \textbf{unidade de tradução} (TU) é o que o
compilador enxerga de cada vez: um \texttt{.cpp} depois de o pré-processador ter
colado dentro dele, textualmente, todos os \texttt{\#include}. É a unidade real
de compilação --- um arquivo-objeto por TU, e nenhuma sabe das outras. Uma
\textbf{fronteira de biblioteca} é o ponto em que código escrito por uma pessoa
passa um valor a código escrito por outra, sem que as duas possam combinar
convenções.

\begin{porque}[o problema: convenções que ninguém verifica]
Antes desses tipos, cada biblioteca inventava a sua forma de dizer as mesmas
quatro ou cinco coisas, e todas tinham o mesmo defeito: nada no tipo obrigava o
outro lado a respeitar a convenção. Para ``pode não haver valor'' usava-se um
ponteiro nulo, um índice $-1$ ou um \texttt{bool} de saída --- e nada exigia a
verificação. Para ``é um destes'', uma \texttt{union} com etiqueta ao lado ---
e nada impedia ler o membro errado. Para ``é qualquer coisa'', \texttt{void*}
--- que perde o tipo por completo. Para ``é uma sequência contígua'', o par
(ponteiro, tamanho) --- dois parâmetros que precisam concordar e que nada obriga
a concordar. O ganho dos tipos de vocabulário não é elegância: é que o erro
passa a ser detectável.
\end{porque}

Cada um dos cinco responde a uma pergunta diferente sobre ``que valor está
aqui?'' (\cref{tab:vocab-perguntas}). A coluna da direita organiza a seção: ela
diz \emph{quando} o conjunto de tipos possíveis é conhecido.

\begin{table}[htbp]
\centering
\caption{As cinco perguntas do vocabulário compartilhado.}
\label{tab:vocab-perguntas}
\footnotesize
\begin{tabular}{@{}lll@{}}
\toprule
\textbf{Tipo} & \textbf{Responde} & \textbf{Tipos conhecidos em} \\
\midrule
\texttt{optional<T>}      & \emph{pode não haver valor}       & compilação, 1 tipo \\
\texttt{variant<Ts...>}   & \emph{é um destes}                & compilação, $N$ tipos \\
\texttt{any}              & \emph{é qualquer coisa}           & \textbf{execução} \\
\texttt{tuple<Ts...>}     & \emph{são todos estes}            & compilação, $N$ tipos \\
\texttt{span<T>}          & \emph{é uma janela sobre aqueles} & compilação, 1 tipo \\
\bottomrule
\end{tabular}
\end{table}

O subprojeto compila com C++23, e não com o C++20 da raiz, porque as operações
monádicas de \texttt{optional} só chegaram nesse padrão (\cref{sec:ambiente}).

% -----------------------------------------------------------------------------
\subsubsection{\classe{std::optional}: a ausência dentro do tipo}
% -----------------------------------------------------------------------------

Uma função que procura pode não achar, e a questão é como ela comunica isso. A
solução tradicional escolhe um valor de retorno para significar ``não achei'':
\texttt{nullptr}, $-1$, \emph{string} vazia. Todas funcionam e todas têm o mesmo
problema --- o valor especial \emph{cabe} no tipo de retorno, então quem chama
pode usá-lo como resultado legítimo, e o compilador não tem o que reclamar.

\classe{std::optional} põe a ausência \textbf{no tipo}. Um
\texttt{std::optional<int>} não é um \texttt{int}: chegar ao \texttt{int} exige
um passo explícito, e é nesse passo que a decisão sobre a ausência é tomada.

\begin{lstlisting}[style=bt, language=C++,
    caption={A ausência declarada na assinatura, em
             \texttt{src/vocabulary\_types/demo\_optional.hpp}.}]
std::optional<int> buscar_idade(const std::string& nome) {
    static const std::unordered_map<std::string, int> pessoas{
        {"Ada", 36}, {"Grace", 45}, {"Alan", 41},
    };
    auto it = pessoas.find(nome);
    if (it == pessoas.end()) return std::nullopt;   // ausência explícita
    return it->second;              // conversão implícita para optional
}
\end{lstlisting}

O \texttt{std::nullopt} é um objeto de tipo próprio que significa ``vazio''; o
\texttt{return it->second;} funciona porque há construtor implícito a partir de
\texttt{T}, de modo que escrever o caso de sucesso não custa nada a mais.

As três formas de acessar diferem no que fazem quando não há valor
(\cref{tab:optional-formas}), e a primeira é a perigosa justamente porque não
falha: falha em silêncio.

\begin{table}[htbp]
\centering
\caption{Extraindo o valor de um \texttt{optional}.}
\label{tab:optional-formas}
\footnotesize
\begin{tabular}{@{}lll@{}}
\toprule
\textbf{Forma} & \textbf{Se estiver vazio} & \textbf{Quando usar} \\
\midrule
\texttt{*opt}             & \textbf{\gls{ub}}                    & já testou com \texttt{if} \\
\texttt{opt.value()}      & lança \texttt{bad\_optional\_access} & quer rede de segurança \\
\texttt{opt.value\_or(x)} & devolve \texttt{x}                   & tem um padrão óbvio \\
\bottomrule
\end{tabular}
\end{table}

\begin{armadilha}[\texttt{*opt} vazio não é exceção nem \emph{crash}]
É \gls{ub}: o padrão não diz o que acontece, e o que acontece na prática é ler
os bytes do \emph{buffer} interno como se fossem um \texttt{T}. Isso costuma
produzir um valor plausível, e o programa segue com um número inventado. Sem
aviso, sem exceção, sem \emph{segmentation fault} garantido. É \textbf{lixo
silencioso}, a pior categoria de erro que existe.
\end{armadilha}

O nome sugere ponteiro, e a sugestão é falsa: não há alocação. O \texttt{T} vive
dentro do próprio \texttt{optional}, num \emph{buffer} alinhado, ao lado de um
\texttt{bool}. O subprojeto prova imprimindo os dois endereços:

\begin{lstlisting}[style=bt, language=bash,
    caption={Saída de \texttt{src/vocabulary\_types/}: o \texttt{T} vive dentro
             do \texttt{optional}.}]
sizeof(int)              =  4
sizeof(optional<int>)    =  8      ← int + flag + padding
&optional  = 0x7ffe9c642e00
&*optional = 0x7ffe9c642e00        ← o mesmo endereço
\end{lstlisting}

\emph{Padding}, nessa saída, são bytes de enchimento que o compilador insere
para que cada membro caia num endereço múltiplo do seu alinhamento: 4 mais 1
dariam 5, mas o objeto precisa ter tamanho múltiplo de 4 para que arranjos dele
fiquem alinhados.

\paragraph{As operações monádicas (C++23).} Encadear operações que podem falhar
gera, na forma ingênua, um \texttt{if} de propagação por passo. As três
operações monádicas eliminam esses \texttt{if} --- se o \texttt{optional}
estiver vazio, a função nem é chamada, e a ausência propaga sozinha:

\begin{lstlisting}[style=bt, language=C++,
    caption={O encadeamento monádico, em
             \texttt{src/vocabulary\_types/demo\_optional.hpp}.}]
nome_do_id(id)                                     // optional<string>
    .and_then(comprimento_se_curto)                // T → optional<U>
    .transform([](int n) { return n * 100; })      // T → U
    .or_else([] { return std::optional<int>{-1}; });
\end{lstlisting}

A distinção entre \texttt{and\_then} e \texttt{transform} é o retorno da função
passada: \texttt{transform} devolve \texttt{U} e a biblioteca embrulha;
\texttt{and\_then} já devolve \texttt{optional<U>} e a biblioteca não embrulha
de novo --- é o que evita o \texttt{optional<optional<U>>}.

\begin{porque}[por que \texttt{optional<T\&>} não existe]
A pergunta que decide é: o que \texttt{opt = outro} deveria fazer com uma
referência dentro? \emph{Religar} a referência para outro objeto, ou
\emph{atribuir} ao objeto referido? As duas respostas são defensáveis, nenhuma é
obviamente correta, e o comitê preferiu proibir a escolher. Na prática o C++ já
tinha o ``\emph{optional reference}'': é o ponteiro cru --- um
\texttt{const Config*} devolvido por uma busca significa exatamente ``opcional e
não-proprietário'', com a desvantagem de não dizer isso no nome. Quando a
assinatura precisa ser inequívoca,
\texttt{optional<std::reference\_wrapper<const T>>} compila, ao custo de um
\texttt{.get()} no uso.
\end{porque}

\begin{armadilha}[\texttt{optional<bool>} tem três estados e o \texttt{if} lê o errado]
O \texttt{operator bool} responde ``\emph{tem valor?}'', não ``\emph{o valor é
verdadeiro?}''. Com um \texttt{bool} dentro, as duas perguntas se confundem:
\begin{lstlisting}[style=bt, language=bash, frame=none,
                   xleftmargin=0pt, xrightmargin=0pt]
vazio        if(o) = false   value_or(false) = false
false        if(o) = true    value_or(false) = false    ← a armadilha
true         if(o) = true    value_or(false) = true
\end{lstlisting}
A linha do meio pega todo mundo: \texttt{if (o)} é verdadeiro porque \emph{há}
valor, ainda que o valor seja \texttt{false}. E \texttt{value\_or(false)}
responde \texttt{false} tanto para o vazio quanto para o \texttt{false}
armazenado --- apagando exatamente a distinção que motivou usar
\texttt{optional}. A saída é modelar com três estados nomeados:
\texttt{enum class Resposta \{ Sim, Nao, NaoInformado \};}. Se o
\texttt{optional<bool>} for inevitável, escreva sempre
\texttt{if (o.has\_value() \&\& *o)}.
\end{armadilha}

\begin{regra}[\texttt{optional} diz ``pode não haver''; não diz ``por quê'']
Use como retorno de busca, \emph{parse} ou consulta; como membro genuinamente
opcional; e para inicialização tardia de membro sem construtor padrão. Não o use
quando o \textbf{motivo} da falha importa --- aí a ferramenta é
\classe{std::expected} (\cref{sec:errors}), que carrega o erro. E não use
\texttt{optional<unique\_ptr<T>>}: o ponteiro já tem estado nulo.
\end{regra}

% -----------------------------------------------------------------------------
\subsubsection{\classe{std::variant}: a soma de tipos que faltava}
% -----------------------------------------------------------------------------

Uma \texttt{struct} é um \emph{produto} de tipos: contém um \texttt{int}
\textbf{e} um \texttt{double} \textbf{e} uma \texttt{string}. A operação
complementar --- \textbf{ou} um \textbf{ou} outro --- chama-se soma de tipos, e
o C++ só tinha uma versão insegura dela: a \texttt{union}, que guarda vários
membros no mesmo espaço mas não guarda \emph{qual} está ativo.

\begin{lstlisting}[style=bt, language=bash,
    caption={Saída de \texttt{src/vocabulary\_types/}: a \texttt{union} não sabe
             quem está ativo; o \texttt{variant} sabe.}]
union: escrevi u.i = 42
  ler u.i → 42 ✓
  ler u.d → seria UNDEFINED BEHAVIOR (membro inativo)
  e NADA no código impede

variant: v = 42
  v.index()                 = 0  (0 = int)
  holds_alternative<double> = false
  std::get<double>(v)       → bad_variant_access
\end{lstlisting}

\classe{std::variant} guarda o mesmo espaço \textbf{mais} o índice do tipo
ativo, e verifica todo acesso: paga um inteiro pequeno e devolve segurança de
tipo. O que ele não paga é alocação --- reserva inline o espaço do maior membro.

Das quatro formas de acessar (\cref{tab:variant-formas}), uma se destaca.
\textbf{Exaustividade} aqui significa que o compilador garante que todos os
tipos possíveis foram tratados: \texttt{std::visit} exige um invocável que
aceite \emph{todos} os tipos, e se faltar um a chamada não compila. É a mesma
propriedade que faz valer a pena um \texttt{switch} sobre \texttt{enum class}
sem \texttt{default} (\cref{sec:syntax}) --- acrescentar um caso quebra o
\emph{build} nos lugares que precisam ser atualizados.

\begin{table}[htbp]
\centering
\caption{Acessando um \texttt{variant}. Só \texttt{visit} verifica exaustividade.}
\label{tab:variant-formas}
\footnotesize
\begin{tabular}{@{}lll@{}}
\toprule
\textbf{Forma} & \textbf{Se o tipo estiver errado} & \textbf{Exaustividade} \\
\midrule
\texttt{get\_if<T>(\&v)} & devolve \texttt{nullptr}      & não verifica \\
\texttt{get<T>(v)}       & \texttt{bad\_variant\_access} & não verifica \\
\texttt{get<I>(v)}       & \texttt{bad\_variant\_access} & não verifica \\
\texttt{visit(f, v)}     & não acontece                  & \textbf{garantida} \\
\bottomrule
\end{tabular}
\end{table}

\paragraph{O \emph{overload idiom}.} O obstáculo prático do \texttt{visit} é que
ele quer \emph{um} objeto invocável que aceite todos os tipos. Um \emph{lambda}
genérico \texttt{[](const auto\& x)} atende, mas o corpo passa a ser o mesmo
para todos --- e o motivo de usar \texttt{variant} era tratar cada um de um
jeito. O \emph{overload idiom} são duas linhas de biblioteca que resolvem isso:

\begin{lstlisting}[style=bt, language=C++,
    caption={O \emph{overload idiom}, em
             \texttt{src/vocabulary\_types/demo\_variant.hpp}.}]
// Herda de todos os lambdas e traz o operator() de cada um para o escopo.
template <typename... Ts>
struct Overload : Ts... { using Ts::operator()...; };

// Deduction guide — em C++20 é opcional, mas explícito não custa.
template <typename... Ts> Overload(Ts...) -> Overload<Ts...>;
\end{lstlisting}

Cada pedaço é um recurso distinto. \texttt{typename... Ts} é um \emph{parameter
pack}: uma lista de tipos de tamanho arbitrário. \texttt{: Ts...} expande essa
lista na cláusula de herança --- a estrutura herda de \emph{cada} tipo, e como
cada \emph{lambda} é uma classe com um \texttt{operator()}, herdar de $N$
\emph{lambdas} traz $N$ funções de chamada. \texttt{using Ts::operator()...;}
expande uma declaração \texttt{using} por tipo, o que é necessário porque, sem
ela, os \texttt{operator()} herdados de bases diferentes não formam um conjunto
de sobrecargas --- ficam escondidos uns pelos outros. Feito isso, a resolução de
sobrecarga normal faz o despacho.

\begin{lstlisting}[style=bt, language=C++,
    caption={Um tratamento por tipo, com exaustividade verificada.}]
using Valor = std::variant<int, double, std::string, bool>;

std::string descrever(const Valor& v) {
    return std::visit(Overload{
        [](int i)                { return fmt::format("inteiro {}", i); },
        [](double d)             { return fmt::format("real {:.2f}", d); },
        [](const std::string& s) { return fmt::format("texto '{}'", s); },
        [](bool b)               { return fmt::format("booleano {}", b); },
    }, v);
}
\end{lstlisting}

Acrescentar um tipo a \texttt{Valor} sem acrescentar o \emph{lambda}
correspondente faz \texttt{descrever} deixar de compilar --- é a exaustividade
em ação. E \texttt{visit} aceita \emph{vários} \texttt{variant} de uma vez, com
o compilador gerando a matriz de combinações: o subprojeto usa
pedra-papel-tesoura, com seis \emph{lambdas} de dois parâmetros mais um
\texttt{[](auto, auto)} que absorve os empates. Fazer o mesmo com herança é o
padrão \emph{Visitor} do GoF (\cref{sec:patterns}) --- duas hierarquias, um
\texttt{accept()} por classe e despacho duplo escrito à mão. Aqui não há
hierarquia, \texttt{accept()} nem \gls{vtable}.

\begin{armadilha}[\texttt{valueless\_by\_exception}: o único estado inválido]
Se a construção do novo valor lançar \emph{durante} uma atribuição, o antigo já
foi destruído e o novo não chegou a existir. O \texttt{variant} fica sem tipo
ativo, e \texttt{index()} devolve \texttt{variant\_npos}:
\begin{lstlisting}[style=bt, language=bash, frame=none,
                   xleftmargin=0pt, xrightmargin=0pt]
antes:  index=0 valueless=false
exceção durante a atribuição: falhei na cópia
depois: index=-1 valueless=true
\end{lstlisting}
Um \texttt{visit} sobre esse estado lança \texttt{bad\_variant\_access}.
Evita-se com três hábitos: pôr um tipo trivial --- ou \classe{std::monostate}
--- como primeira alternativa; preferir \texttt{emplace} à atribuição; e marcar
o construtor de movimento como \texttt{noexcept}. É a Rule of 5 do
\cref{sec:value-semantics} cobrando de novo.
\end{armadilha}

\classe{std::monostate} é um tipo vazio cuja única função é servir de primeira
alternativa, e ele resolve um problema concreto: um \texttt{variant} só tem
construtor padrão se o \emph{primeiro} tipo da lista tiver.
\texttt{variant<SemDefault, int> v;} não compila;
\texttt{variant<std::monostate, SemDefault, int> v;} compila e nasce no estado
``nada ainda''.

\paragraph{Um \texttt{variant} recursivo.} Um \texttt{variant} não pode conter a
si mesmo diretamente --- o tamanho seria infinito, já que reserva espaço inline
para o maior tipo. A saída é a indireção:

\begin{lstlisting}[style=bt, language=C++,
    caption={Uma árvore de expressão sem uma única função virtual, em
             \texttt{src/vocabulary\_types/demo\_variant.hpp}.}]
using Expressao = std::variant<double,
                               std::unique_ptr<Soma>,
                               std::unique_ptr<Multiplicacao>>;

struct Soma          { Expressao esq, dir; };
struct Multiplicacao { Expressao esq, dir; };

double avaliar(const Expressao& e) {
    return std::visit(Overload{
        [](double d) { return d; },
        [](const std::unique_ptr<Soma>& s) {
            return avaliar(s->esq) + avaliar(s->dir);
        },
        [](const std::unique_ptr<Multiplicacao>& m) {
            return avaliar(m->esq) * avaliar(m->dir);
        },
    }, e);
}
\end{lstlisting}

Nenhuma classe base, nenhum método virtual, nenhum \texttt{accept()}. A
hierarquia é \textbf{fechada} --- e é por ser fechada que o compilador consegue
verificar a exaustividade.

\paragraph{O \emph{expression problem}.} A comparação com herança não tem
vencedor absoluto, e isso tem nome: os dois eixos de extensão de um programa ---
acrescentar \emph{tipos} e acrescentar \emph{operações} --- puxam em direções
opostas, e a estrutura que facilita um dificulta o outro
(\cref{tab:variant-vs-heranca}). Com \texttt{variant}, uma operação nova é uma
função nova e nenhum tipo existente muda; mas um tipo novo obriga a revisitar
todos os \texttt{visit}. Com herança é o inverso: um tipo novo pode nascer em
outra biblioteca, e uma operação nova obriga a mexer na base e em todas as
derivadas. Escolha pelo eixo que vai variar mais.

\begin{table}[htbp]
\centering
\caption{\texttt{variant} favorece operações novas; herança favorece tipos novos.}
\label{tab:variant-vs-heranca}
\footnotesize
\begin{tabular}{@{}lcc@{}}
\toprule
\textbf{Critério} & \texttt{variant} & \textbf{herança virtual} \\
\midrule
tipos novos \emph{externos} & não          & \textbf{sim} \\
operações novas             & \textbf{sim} & recompila tudo \\
exaustividade verificada    & \textbf{sim} & não \\
alocação por nó             & opcional     & sempre \\
\bottomrule
\end{tabular}
\end{table}

% -----------------------------------------------------------------------------
\subsubsection{\classe{std::any} e o apagamento de tipo}
% -----------------------------------------------------------------------------

\texttt{optional<T>} sabe o tipo; \texttt{variant<Ts...>} sabe o \emph{conjunto}
de tipos; \classe{std::any} não sabe nada em compilação: aceita qualquer tipo
copiável e só descobre qual é em execução.

O mecanismo chama-se \gls{typeerasure}: guardar, junto com o objeto, ponteiros
para as funções que sabem operar sobre ele --- copiar, destruir, informar o
\texttt{typeid} --- de modo que o código que manipula o \texttt{any} faça o
necessário sem conhecer o tipo. É o princípio do \gls{vtable}, com a diferença
de que a tabela não vive dentro do objeto e o tipo armazenado não precisa herdar
de nada. A recuperação é feita por \texttt{any\_cast}, em duas formas: a que
recebe o objeto e lança \texttt{bad\_any\_cast}, e a que recebe um
\emph{ponteiro} e devolve \texttt{nullptr}. O casamento de tipo é
\textbf{exato} --- um \texttt{any} com \texttt{double} não responde a
\texttt{any\_cast<int>}.

\paragraph{O custo, contado.} Em vez de afirmar que \texttt{any} aloca, o
subprojeto \emph{substitui o \texttt{operator new} global} por uma versão que
conta chamadas e bytes, e mede:

\begin{lstlisting}[style=bt, language=bash,
    caption={Saída de \texttt{src/vocabulary\_types/}: alocações contadas ao
             construir um objeto.}]
sizeof(std::any)                     = 16
any<Pequeno> ( 4B) → 0 alocação(ões)
any<Medio>   (16B) → 1 alocação(ões), 16 bytes
any<Grande>  (64B) → 1 alocação(ões), 64 bytes
variant<...Grande>  → 0 alocação(ões)      ← nunca aloca
\end{lstlisting}

A primeira linha é a \textbf{otimização de objeto pequeno} (SBO): a
implementação reserva alguns bytes dentro do próprio objeto e, se o valor
couber, guarda-o ali sem \emph{heap}. Na libstdc++ o limite do SBO do
\texttt{any} é \texttt{sizeof(void*)} --- oito bytes numa máquina de 64 bits.
Acima disso, \emph{cada} \texttt{any} construído é um \texttt{malloc} mais um
\texttt{free}. A última linha é o contraste que importa: o \texttt{variant}
\textbf{nunca} aloca. Essa diferença estrutural pesa muito mais que a velocidade
bruta de uma chamada isolada --- num laço quente, o que se mede é a pressão
sobre o alocador.

\begin{regra}[a pergunta que elimina o \texttt{any} quase sempre]
Pergunte: \emph{``eu consigo listar os tipos possíveis?''} \textbf{Sim}
$\rightarrow$ \texttt{variant}, que ganha exaustividade e zero alocação.
\textbf{Não, mas todos fazem a mesma coisa} $\rightarrow$ interface virtual: os
tipos são abertos e o que os une é o comportamento. \textbf{Não, e cada um faz
coisa diferente} $\rightarrow$ aí sim, \texttt{any}. O sintoma de \texttt{any}
mal usado é uma cascata de \texttt{any\_cast} encadeados --- e essa cascata
\emph{é} um \texttt{variant}, escrito à mão e sem exaustividade.
\end{regra}

Onde \texttt{any} é a resposta certa: quando o conjunto de tipos é aberto
\textbf{por projeto}, definido por quem \emph{usa} a biblioteca. O subprojeto
mostra o caso canônico --- uma classe \texttt{Propriedades} sobre um
\texttt{std::map<std::string, std::any>} cujo método de leitura devolve
\texttt{std::optional<T>}, combinando ``pode não existir'' e ``pode ser de
qualquer tipo''. É assim que funcionam os ``\emph{user data}'' de bibliotecas de
interface gráfica e os \emph{callbacks} de APIs em C.

% -----------------------------------------------------------------------------
\subsubsection{\texttt{tuple}, \texttt{pair} e os \emph{structured bindings}}
% -----------------------------------------------------------------------------

Uma \texttt{struct} nomeia os campos; uma \classe{std::tuple} não. É a única
diferença relevante, e ela decide quando usar cada uma.

O problema histórico da \texttt{tuple} era o acesso: \texttt{std::get<0>(r)} não
diz nada sobre o que é o campo zero. Os \emph{structured bindings} do C++17
mudaram essa conta ao dar nomes no ponto de uso --- e funcionam sobre
\texttt{tuple}, \texttt{pair}, arranjo e \texttt{struct}:

\begin{lstlisting}[style=bt, language=C++,
    caption={\emph{Structured bindings} nos quatro contextos mais comuns.}]
auto [nome, idade, salario] = obter_registro();   // tuple
auto [x, y] = ponto;                              // struct, sem métodos

for (const auto& [fruta, qtd] : estoque)          // map: chave e valor
    fmt::print("{} → {}\n", fruta, qtd);

if (auto [it, inserido] = m.try_emplace("nova", 1); inserido)
    fmt::print("inseriu\n");                      // pair<iterator,bool>
\end{lstlisting}

\begin{armadilha}[os nomes do \emph{binding} são novos, não \emph{aliases}]
\texttt{auto [a,b]} \textbf{copia} o objeto inteiro antes de dar nome às partes;
\texttt{auto\& [a,b]} referencia; \texttt{const auto\& [a,b]} é o padrão seguro.
Escrever \texttt{auto} por reflexo sobre um objeto grande copia o objeto sem
avisar --- e alterar \texttt{a} não altera o original, o que costuma ser
descoberto tarde.
\end{armadilha}

\paragraph{\texttt{std::tie}.} Constrói uma \texttt{tuple} de \emph{referências}
aos argumentos, sem copiar. Como \texttt{tuple} já tem \texttt{operator<}
lexicográfico, isso resolve comparação por múltiplos critérios numa linha:

\begin{lstlisting}[style=bt, language=C++,
    caption={Três usos de \texttt{std::tie}, em
             \texttt{src/vocabulary\_types/demo\_tuple.hpp}.}]
return std::tie(a.idade, a.nome) < std::tie(b.idade, b.nome);  // desempate
std::tie(nome, idade, salario) = obter_registro();     // variáveis existentes
std::tie(std::ignore, so_a_idade, std::ignore) = obter_registro();
\end{lstlisting}

A alternativa escrita à mão são duas linhas com um \texttt{if}; com quatro
critérios, viram sete linhas e um defeito esperando para acontecer.

\begin{armadilha}[\texttt{tie} guarda referências: nunca o devolva]
O resultado referencia os objetos passados. Devolvê-lo de uma função deixa
referências penduradas para locais que já morreram. Para devolver,
\texttt{std::make\_tuple}, que copia.
\end{armadilha}

\texttt{std::apply} desempacota uma \texttt{tuple} nos parâmetros de uma função
e \texttt{std::make\_from\_tuple} faz o mesmo para um construtor: os dois são a
base de qualquer mecanismo que guarde argumentos agora e chame depois ---
\texttt{std::bind}, filas de \emph{thread pool}, sistemas de tarefas
(\cref{sec:concurrency}).

Percorrer uma \texttt{tuple} não é possível com \emph{range-for}, porque cada
elemento tem um \textbf{tipo} diferente e não há \texttt{operator++} que faça
sentido. Exige desdobrar em compilação, o que uma \emph{fold expression} resolve
em uma linha: \texttt{std::apply([\&f](const auto\&... e) \{ (f(e), ...); \},
t);}. Uma \emph{fold expression} (C++17) aplica um operador binário a todos os
elementos de um \emph{pack}; aqui o operador é a vírgula, e \texttt{(f(e), ...)}
expande em \texttt{f(e1), f(e2), f(e3)} --- o corpo é instanciado uma vez por
elemento, com o tipo certo.

Quanto a escolher (\cref{tab:tuple-ou-struct}): o \texttt{sizeof} é
\textbf{idêntico} --- o subprojeto verifica ---, então a decisão é inteiramente
de legibilidade. E desde o C++20 a inicialização designada
(\texttt{return R\{.encontrou = true, .indice = 3\};}) tirou o último argumento
prático a favor da \texttt{tuple} em código de aplicação.

\begin{table}[htbp]
\centering
\caption{A escolha é 100\% de legibilidade: o custo é o mesmo.}
\label{tab:tuple-ou-struct}
\footnotesize
\begin{tabular}{@{}ll@{}}
\toprule
\textbf{Use \texttt{tuple}/\texttt{pair} quando\ldots} & \textbf{Use \texttt{struct} quando\ldots} \\
\midrule
os campos não têm nome natural         & o agregado é do \textbf{domínio} \\
é interno e efêmero                    & atravessa uma API pública \\
precisa de \texttt{tie}/\texttt{apply} & tem invariantes a manter \\
é código genérico sobre $N$ tipos      & tem métodos próprios \\
são 2 valores óbvios (\texttt{pair})   & tem 3 ou mais campos \\
\bottomrule
\end{tabular}
\end{table}

% -----------------------------------------------------------------------------
\subsubsection{\classe{std::span}: a janela sobre memória alheia}
% -----------------------------------------------------------------------------

A assinatura mais comum do C legado é o par (ponteiro, tamanho):

\begin{lstlisting}[style=bt, language=C++,
    caption={O parâmetro que não pode discordar de si mesmo, em
             \texttt{src/vocabulary\_types/demo\_span.hpp}.}]
int soma_antiga(const int* dados, std::size_t n);  // soma_antiga(bruto, 500) compila
int soma_span(std::span<const int> dados);          // o tamanho vem do objeto
\end{lstlisting}

A primeira funciona e é rápida, mas separa duas informações que só fazem sentido
juntas. \texttt{soma\_antiga(v, 500)} sobre um arranjo de cinco elementos
compila sem um único aviso e lê 495 inteiros de memória que não é dela. E o
corpo da função não pode usar \emph{range-for} nem os algoritmos da biblioteca,
porque um ponteiro cru não é um \emph{range}.

\classe{std::span} é esse par unificado num tipo que \emph{é} um \emph{range}.
Não possui os dados, aceita arranjo C, \classe{std::array},
\classe{std::vector} e qualquer sequência contígua, e o tamanho vem sempre do
objeto de origem --- é impossível errar.

O segundo parâmetro é o \emph{extent}: o número de elementos, quando conhecido
em compilação. A diferença aparece no \texttt{sizeof} medido:

\begin{lstlisting}[style=bt, language=bash,
    caption={Saída de \texttt{src/vocabulary\_types/}: o tamanho no tipo não
             ocupa espaço no objeto.}]
sizeof(span<int>)    = 16      ponteiro + tamanho
sizeof(span<int, 3>) =  8      só o ponteiro — o 3 está no TIPO
sizeof(vector<int>)  = 24      ponteiro + size + capacity
\end{lstlisting}

Um \texttt{span<double, 3>} pode declarar \texttt{static\_assert(v.extent ==
3);} dentro da própria função. A conversão de estático para dinâmico é
implícita; a inversa exige \texttt{.first<3>()}, e quem escreve assume o risco.
As janelas --- \texttt{first(n)}, \texttt{last(n)}, \texttt{subspan(i, n)} ---
devolvem outro \texttt{span} sobre o mesmo \emph{buffer}, sem copiar nada.

\begin{armadilha}[\texttt{subspan} fora dos limites é \gls{ub}, não exceção]
As versões com \emph{template} (\texttt{.first<4>()}) verificam em compilação
\emph{quando o extent é estático}; as versões com argumento de execução não
verificam nada. Prefira as primeiras sempre que o tamanho for conhecido.
\end{armadilha}

Há duas coisas diferentes que podem ser congeladas, e a sintaxe engana
(\cref{tab:span-constancia}). \texttt{span<const T>} protege os
\emph{elementos}: não dá para escrever pela janela. \texttt{const span<T>}
protege a \emph{janela}: não dá para reapontar o \texttt{span}, mas escrever nos
elementos compila normalmente. Numa assinatura,
\texttt{std::span<const T>} é o padrão para ``só lê''.

\begin{table}[htbp]
\centering
\caption{\texttt{span<const T>} contra \texttt{const span<T>}.}
\label{tab:span-constancia}
\footnotesize
\begin{tabular}{@{}ll@{}}
\toprule
\textbf{Quer proteger} & \textbf{Escreva} \\
\midrule
os \textbf{elementos} & \texttt{span<const T>} --- 95\% dos casos \\
a \textbf{janela}     & \texttt{const span<T>} \\
ambos                 & \texttt{const span<const T>} \\
\bottomrule
\end{tabular}
\end{table}

Um \emph{borrowed range} é um \emph{range} cujos iteradores continuam válidos
depois de o objeto \emph{range} morrer. Vale para \texttt{span} e
\classe{std::string\_view} porque eles não possuem nada:

\begin{lstlisting}[style=bt, language=bash,
    caption={Saída de \texttt{src/vocabulary\_types/}:
             \texttt{std::ranges::borrowed\_range}.}]
span<int>          : true
string_view        : true
vector<int>        : false     ← o objeto POSSUI os dados
vector<int>&       : true      ← a referência não possui
\end{lstlisting}

A consequência prática é que os algoritmos de \cabecalho{ranges} aceitam um
\texttt{span} temporário e devolvem iterador utilizável, mas recusam um
\texttt{vector} temporário --- devolvem \texttt{std::ranges::dangling}
(\cref{sec:ranges}).

\begin{armadilha}[um \texttt{span} é tão frágil quanto um par de iteradores]
\begin{lstlisting}[style=bt, language=C++, frame=none,
                   xleftmargin=0pt, xrightmargin=0pt]
std::span<int> s{v};
v.push_back(4);      // realocou → s está pendurado
s[0] = 1;            // heap-use-after-free
\end{lstlisting}
Porque \emph{é} isso: um ponteiro e um tamanho. Vale a mesma tabela de
invalidação do \cref{sec:stl}. As três regras de uso: (1)~\texttt{span} é
\textbf{parâmetro} --- o uso dominante e sempre seguro; (2)~como retorno, só se
o \emph{buffer} sobreviver à chamada, o que na prática significa apontar para um
membro do objeto; (3)~como membro, exige documentar quem é o dono. O subprojeto
recomenda \texttt{-Db\_sanitize=address,undefined} justamente por isto.
\end{armadilha}

\begin{regra}[o resumo do vocabulário]
Função que pode não achar nada: \texttt{optional<T>}; e se o \emph{motivo}
importa, \texttt{expected<T,E>}. ``Um destes $N$ tipos'', lista fechada:
\texttt{variant} com \texttt{visit} e o \emph{overload idiom}. Tipos definidos
por quem \textbf{usa} a API: \texttt{any} --- e, se aparecer cascata de
\texttt{any\_cast}, era \texttt{variant}. Dois valores sem nome natural:
\texttt{pair}; três ou mais, ou API pública: \texttt{struct} nomeado. Parâmetro
que recebe sequência contígua: \texttt{span<const T>}.
\end{regra}

% =============================================================================
\subsection{Computação em tempo de compilação}
\label{sec:ct-fundo}
% =============================================================================

A segunda metade do capítulo troca de assunto sem trocar de estratégia. Se pôr
informação no tipo permite ao compilador verificar, pôr \emph{cálculo} no tempo
de compilação permite a ele \textbf{executar} --- e, o que é mais interessante,
\emph{provar}.

\begin{porque}[o argumento não é desempenho: é a ausência de \gls{ub}]
O avaliador de expressões constantes é obrigado, pelo padrão, a recusar qualquer
operação cujo comportamento não esteja definido. Isso transforma uma classe
inteira de defeitos de execução em erros de compilação:
\begin{lstlisting}[style=bt, language=C++, frame=none,
                   xleftmargin=0pt, xrightmargin=0pt]
constexpr int x = 1 / 0;            // erro de compilação, não crash
constexpr int y = arr[10];          // erro se estiver fora dos limites
constexpr int z = *(int*)nullptr;   // erro
\end{lstlisting}
Código que passa sob \texttt{static\_assert} é código \textbf{provadamente}
livre de \gls{ub} naquele caminho. Nenhuma bateria de testes dá essa garantia:
testes cobrem as entradas que alguém pensou em escrever; isto cobre o caminho
inteiro. A vantagem de desempenho é real, mas é a segunda melhor razão.
\end{porque}

% -----------------------------------------------------------------------------
\subsubsection{\texttt{constexpr} é permissão, não garantia}
% -----------------------------------------------------------------------------

O erro conceitual mais comum é ler \texttt{constexpr} como ``isto sempre roda na
compilação''. É uma \textbf{permissão}: ``isto \emph{pode} rodar na compilação,
se os argumentos permitirem''. A mesma função serve aos dois mundos, e quem
decide é o contexto da chamada:

\begin{lstlisting}[style=bt, language=C++,
    caption={A mesma função nos dois mundos, em
             \texttt{src/compile\_time/demo\_constexpr.hpp}.}]
constexpr int fatorial(int n) {
    int r = 1;
    for (int i = 2; i <= n; ++i) r *= i;   // laço: legal desde C++14
    return r;
}

constexpr int em_compilacao = fatorial(10);   // avaliado AGORA
int n = 5;
const int em_runtime = fatorial(n);           // vira código de máquina
\end{lstlisting}

Para \textbf{forçar} a avaliação existem três contextos, do mais fraco para o
mais forte: \texttt{constexpr int a = f();} força naquela declaração;
\texttt{static\_assert(f() == x);} força e ainda verifica o valor; e
\texttt{std::array<int, f()>} força, porque um parâmetro de \emph{template} que
é valor é o contexto mais exigente que a linguagem tem. Sem nenhum desses, o
compilador \emph{pode} avaliar em compilação --- e com \texttt{-O2} normalmente
avalia ---, mas não é obrigado. Para \emph{garantir} que uma função nunca rode
em execução, a ferramenta é \texttt{consteval}.

O que cabe dentro de uma função \texttt{constexpr} cresceu a cada padrão
(\cref{tab:evolucao-constexpr}). Conhecer essa tabela evita duas frustrações:
tentar escrever um laço em C++11 e concluir que \texttt{constexpr} é inútil, e
tentar devolver um \texttt{vector} em C++20 e concluir que o compilador está com
defeito.

\begin{table}[htbp]
\centering
\caption{O que passou a caber numa função \texttt{constexpr}, por padrão.}
\label{tab:evolucao-constexpr}
\footnotesize
\begin{tabular}{@{}ll@{}}
\toprule
\textbf{Padrão} & \textbf{Passou a aceitar} \\
\midrule
C++11 & praticamente só \texttt{return expr;} --- uma única instrução \\
C++14 & laços, variáveis locais, \texttt{if}/\texttt{switch}, vários \texttt{return} \\
C++17 & \texttt{if constexpr}, \emph{lambdas} \texttt{constexpr} \\
C++20 & alocação transitória, \texttt{try}/\texttt{catch}, \texttt{virtual}, \texttt{vector}, \texttt{string} \\
C++23 & \texttt{if consteval}, \texttt{goto}, variáveis \texttt{static} locais \\
\bottomrule
\end{tabular}
\end{table}

O C++20 permitiu alocar durante uma avaliação em compilação, o que trouxe
\classe{std::vector} e \classe{std::string} para dentro do \texttt{constexpr}. A
permissão vem com uma condição --- \textbf{toda alocação feita durante a
avaliação precisa ser liberada antes de a avaliação terminar} ---, a chamada
alocação \emph{transitória}. Isso permite usar um \texttt{vector} como memória
de rascunho, e o subprojeto resolve o crivo de Eratóstenes inteiro antes de o
programa existir:

\begin{lstlisting}[style=bt, language=C++,
    caption={O crivo de Eratóstenes rodando no compilador.}]
constexpr int soma_dos_primos_ate(int limite) {
    std::vector<bool> composto(limite + 1, false);      // aloca
    int soma = 0;
    for (int i = 2; i <= limite; ++i) {
        if (composto[i]) continue;
        soma += i;
        for (int j = i * i; j <= limite; j += i) composto[j] = true;
    }
    return soma;                       // o vector é destruído aqui
}

constexpr int soma_primos = soma_dos_primos_ate(1000);
static_assert(soma_primos == 76127);
\end{lstlisting}

\begin{armadilha}[\texttt{constexpr std::vector<int> v = \{1,2,3\};} não compila]
A regra transitória proíbe que a alocação sobreviva ao fim da avaliação --- e
uma variável \texttt{constexpr} \emph{é} o resultado que sobrevive. A memória do
compilador não existe mais quando o programa roda; não haveria para onde
apontar. Logo, um \texttt{vector} pode aparecer no \emph{meio} de um cálculo
\texttt{constexpr} e nunca ser o \emph{resultado} dele. Para devolver dados de
uma função \texttt{constexpr}, o tipo é \classe{std::array}, cujo armazenamento
é o próprio objeto.
\end{armadilha}

Um construtor \texttt{constexpr} permite que objetos existam em contexto de
compilação, e desde o C++14 métodos \texttt{constexpr} não são implicitamente
\texttt{const} --- podem modificar o objeto durante a avaliação. O
\texttt{Racional} do subprojeto normaliza \texttt{2/4} para \texttt{1/2} no
construtor, e \texttt{static\_assert(meio.numerador() == 1);} verifica isso sem
executar nada. Onde a mutação é necessária, o idioma é um \emph{lambda} chamado
imediatamente, que cria um escopo mutável e devolve a constante pronta.

Uma avaliação \texttt{constexpr} não pode fazer entrada/saída --- não há sistema
operacional na compilação ---, chamar função não-\texttt{constexpr}, ler
variável não-\texttt{constexpr}, usar \texttt{reinterpret\_cast} nem deixar
alocação viva ao fim. E há um limite prático: o GCC impõe um teto de passos de
avaliação (\texttt{-fconstexpr-ops-limit}, padrão 33.554.432); ultrapassá-lo
produz \emph{``constexpr evaluation operation count exceeded''}.

% -----------------------------------------------------------------------------
\subsubsection{\texttt{consteval} e \texttt{constinit}: as garantias}
% -----------------------------------------------------------------------------

\texttt{constexpr} deixa duas perguntas sem resposta, e o C++20 acrescentou uma
palavra-chave para cada uma (\cref{tab:tres-palavras}). Funções
\texttt{consteval} chamam-se \emph{immediate functions}: toda chamada precisa
ser, ela própria, expressão constante. Chamar com uma variável não compila, com
a mensagem \emph{``call to consteval function is not a constant expression''}.

\begin{table}[htbp]
\centering
\caption{As três palavras e o que cada uma exige.}
\label{tab:tres-palavras}
\footnotesize
\begin{tabular}{@{}lcc@{}}
\toprule
\textbf{Palavra} & \textbf{Avaliação em compilação} & \textbf{Argumento de execução} \\
\midrule
(nenhuma)          & não                  & aceita \\
\texttt{constexpr} & \textbf{permitida}   & aceita \\
\texttt{consteval} & \textbf{obrigatória} & \textbf{erro} \\
\bottomrule
\end{tabular}
\end{table}

A regra que confunde é que \texttt{consteval} é \textbf{contagiosa para cima}:
uma função \texttt{constexpr} \emph{pode} ser chamada em execução; se ela
repassasse o próprio parâmetro a uma \texttt{consteval}, essa chamada teria de
acontecer em execução --- proibido por definição. O compilador recusa já na
definição:

\begin{lstlisting}[style=bt, language=C++,
    caption={Quem repassa um parâmetro para uma \texttt{consteval} precisa ser
             \texttt{consteval}.}]
consteval int dobro_imediato(int x) { return x * 2; }

constexpr int f(int x) { return dobro_imediato(x); }      // ✗ não compila
consteval int g(int x) { return dobro_imediato(x) + 1; }  // ✓ contexto imediato
constexpr int h(int x) { return x + dobro_imediato(10); } // ✓ argumento literal
\end{lstlisting}

A aplicação mais útil é a \textbf{validação de argumento literal}. Como um
\texttt{throw} que escapa de uma avaliação constante é erro de compilação, um
construtor \texttt{consteval} recusa valores inválidos antes de o programa
existir:

\begin{lstlisting}[style=bt, language=C++,
    caption={Validação em compilação, em
             \texttt{src/compile\_time/demo\_consteval.hpp}.}]
class Porta {
public:
    consteval explicit Porta(int valor) : valor_(valor) {
        if (valor < 1 || valor > 65535)
            throw "porta fora do intervalo 1..65535";
    }
    constexpr int valor() const { return valor_; }
private:
    int valor_;
};

constexpr Porta http{80};       // ✓
// constexpr Porta ruim{70000}; // ✗ NÃO COMPILA — a mensagem cita o throw
\end{lstlisting}

É exatamente o mecanismo do \classe{std::format} do C++20: a \emph{string} de
formato é validada por uma \texttt{consteval}, e por isso
\texttt{std::format("\{:d\}", "texto")} falha na \textbf{compilação}, não com
exceção em produção. O subprojeto reproduz a ideia com uma
\texttt{contar\_placeholders} que recusa uma chave sem fechamento.

\paragraph{\texttt{constinit} e o \emph{static initialization order fiasco}.}
Variáveis globais de unidades de tradução diferentes são inicializadas em ordem
\textbf{não especificada} --- ordem que, na prática, depende de como o
\emph{linker} arrumou os objetos. Se a inicialização de uma global em
\texttt{b.cpp} lê uma global de \texttt{a.cpp}, o programa pode ler um objeto
ainda não construído. Esse é o fiasco: um defeito que aparece e desaparece
conforme a ordem de link muda.

\texttt{constinit} exige inicialização \textbf{estática} --- feita na compilação
e gravada direto no binário. Se o valor já está lá quando o programa começa, não
há ordem a respeitar; e o compilador recusa qualquer inicialização que não caiba
na compilação, de modo que o fiasco vira erro de \emph{build}. A terceira linha
da \cref{tab:constinit} é a razão de a palavra existir: sem ela, a única forma
de garantir inicialização estática era tornar a variável \texttt{constexpr}, o
que a congela. A outra saída clássica é o \emph{Meyers singleton}
(\cref{sec:solid}).

\begin{table}[htbp]
\centering
\caption{\texttt{constinit} não implica \texttt{const}.}
\label{tab:constinit}
\footnotesize
\begin{tabular}{@{}lll@{}}
\toprule
\textbf{Palavra} & \textbf{Inicialização} & \textbf{Mutável depois} \\
\midrule
\texttt{const}     & qualquer            & não \\
\texttt{constexpr} & compilação          & não \\
\texttt{constinit} & \textbf{compilação} & \textbf{sim} \\
\bottomrule
\end{tabular}
\end{table}

\paragraph{\texttt{if consteval} (C++23).} Às vezes o algoritmo ótimo difere
entre os dois mundos: em execução convém o intrínseco do processador; em
compilação, ele não existe e é preciso um laço.

\begin{lstlisting}[style=bt, language=C++,
    caption={Uma função, dois algoritmos, um resultado.}]
constexpr int contar_bits(unsigned x) {
    if consteval {
        int n = 0;                          // caminho de COMPILAÇÃO
        while (x) { n += static_cast<int>(x & 1u); x >>= 1; }
        return n;
    } else {
        return __builtin_popcount(x);       // caminho de EXECUÇÃO
    }
}
\end{lstlisting}

\begin{armadilha}[\texttt{std::is\_constant\_evaluated()} tem três pegadinhas]
A função do C++20 que precedeu o \texttt{if consteval} erra de três maneiras.
\textbf{Um:} dentro de um \texttt{if} comum, responde \texttt{true} sempre que a
função \emph{é} \texttt{constexpr}, mesmo quando aquela chamada está sendo
compilada para execução. \textbf{Dois:}
\texttt{if constexpr (std::is\_constant\_evaluated())} é \textbf{sempre}
verdadeiro, porque a própria condição é avaliada como constante --- o ramo
\texttt{else} vira código morto, silenciosamente. \textbf{Três:} o ramo de um
\texttt{if} comum \textbf{não é contexto imediato}, então chamar uma função
\texttt{consteval} lá dentro é erro de compilação, mesmo que o ramo só fosse
executado durante a avaliação constante. O \texttt{if consteval} \emph{cria}
esse contexto, e é essa a diferença que justificou uma palavra-chave nova.
\end{armadilha}

% -----------------------------------------------------------------------------
\subsubsection{\texttt{static\_assert} e \texttt{type\_traits}}
% -----------------------------------------------------------------------------

\texttt{static\_assert} é a única forma de fazer o compilador falhar de
propósito, com uma mensagem sua. Combinado com \cabecalho{type\_traits},
transforma suposições implícitas sobre tipos em contratos verificados. A regra
que orienta o uso: \emph{se você escreveu o comentário ``T precisa ser
copiável'', esse comentário deveria ser um \texttt{static\_assert}}.

\begin{regra}[a mensagem é o produto]
Um \texttt{static\_assert} sem mensagem faz o compilador repetir a
\emph{expressão}, que raramente ajuda. A mensagem deve dizer \textbf{o que
fazer}, não o que falhou --- quem lê já sabe o que falhou.
\begin{lstlisting}[style=bt, language=C++, frame=none,
                   xleftmargin=0pt, xrightmargin=0pt]
static_assert(std::is_trivially_copyable_v<T>,
    "serializar() copia bytes crus e só aceita tipos trivialmente "
    "copiáveis. Tipos com ponteiros, std::string ou destrutor próprio "
    "precisam de uma sobrecarga específica.");
\end{lstlisting}
Compare com a alternativa: \emph{``static assertion failed''}.
\end{regra}

Um \emph{type trait} é um \emph{template} que responde uma pergunta sobre um
tipo em compilação. Os que respondem sim ou não terminam em \texttt{\_v}; os de
transformação terminam em \texttt{\_t} e produzem outro tipo
(\cref{tab:traits}). Entre os de transformação, um explica um comportamento que
todo mundo já encontrou: \texttt{decay\_t<const int\&>} é \texttt{int} e
\texttt{decay\_t<int[5]>} é \texttt{int*} --- exatamente o que \texttt{auto} faz
na dedução por valor.

\begin{table}[htbp]
\centering
\caption{O catálogo de \emph{traits} que aparece de verdade.}
\label{tab:traits}
\footnotesize
\begin{tabular}{@{}ll@{}}
\toprule
\textbf{Pergunta} & \textbf{Trait} \\
\midrule
dá para \texttt{memcpy}?                  & \texttt{is\_trivially\_copyable\_v} \\
tem \gls{vtable}?                         & \texttt{is\_polymorphic\_v} \\
o \emph{move} é barato \textbf{e} seguro? & \texttt{is\_nothrow\_move\_constructible\_v} \\
dá para inicializar com \texttt{\{\}}?    & \texttt{is\_aggregate\_v} \\
herda de X? converte para X?              & \texttt{is\_base\_of\_v}, \texttt{is\_convertible\_v} \\
é chamável com estes argumentos?          & \texttt{is\_invocable\_v} \\
\bottomrule
\end{tabular}
\end{table}

A pergunta ``este tipo tem o método X?'' não tem \emph{trait} pronto. A forma
clássica usa \texttt{std::void\_t} com especialização parcial e \gls{sfinae}; a
moderna usa um \gls{concept}:

\begin{lstlisting}[style=bt, language=C++,
    caption={O mesmo teste, em C++17 e em C++20, em
             \texttt{src/compile\_time/demo\_static\_assert.hpp}.}]
// C++17: void_t + especialização parcial (8 linhas)
template <typename, typename = void>
struct tem_serializar : std::false_type {};

template <typename T>
struct tem_serializar<T,
        std::void_t<decltype(std::declval<const T&>().serializar())>>
    : std::true_type {};

// C++20: um concept diz a mesma coisa em 3 linhas — e diz mais
template <typename T>
concept Serializavel = requires(const T& t) {
    { t.serializar() } -> std::convertible_to<std::string>;
};
\end{lstlisting}

O \emph{concept} ganha em três frentes, e a terceira surpreende: além de ser
mais curto e de dar mensagem de erro na chamada em vez de no meio da
instanciação, ele verifica também o \textbf{tipo de retorno}. Um tipo cujo
\texttt{serializar()} devolve \texttt{int} passa no \emph{detection idiom} ---
que só pergunta se a expressão é válida --- e é corretamente rejeitado pelo
\emph{concept}.

Os dois não são substitutos: falham em \textbf{momentos diferentes}, e isso muda
a resolução de sobrecarga (\cref{tab:assert-concept}).
\texttt{com\_assert(3.14)} é erro --- o \emph{template} foi escolhido e só então
o \emph{assert} falhou, sem chance de outra sobrecarga assumir.
\texttt{com\_concept(3.14)} funciona: a restrição removeu o \emph{template} do
conjunto de candidatos, e \texttt{void com\_concept(double)} venceu. Use
\emph{concept} para \textbf{restringir}; \texttt{static\_assert} para
\textbf{afirmar} fatos sem alternativa.

\begin{table}[htbp]
\centering
\caption{O momento da falha decide se há alternativa.}
\label{tab:assert-concept}
\footnotesize
\begin{tabular}{@{}lll@{}}
\toprule
\textbf{Ferramenta} & \textbf{Quando falha} & \textbf{Afeta a sobrecarga?} \\
\midrule
\texttt{static\_assert} & \textbf{após} a escolha    & não \\
\emph{concept}          & \textbf{durante} a escolha & \textbf{sim} \\
\bottomrule
\end{tabular}
\end{table}

Há um domínio em que \emph{concepts} não substituem nada: \emph{layout} e
\gls{abi}. Estruturas que atravessam uma fronteira binária --- protocolo de
rede, formato de arquivo, chamada para outra linguagem --- têm requisitos de
tamanho e de deslocamento de campo que \emph{concepts} não expressam, porque
restringem \emph{templates} e não dizem nada sobre disposição de bytes.

\begin{lstlisting}[style=bt, language=C++,
    caption={O contrato do protocolo, verificado no \emph{build}.}]
static_assert(sizeof(CabecalhoDeRede) == 8,
    "O cabeçalho do protocolo tem 8 bytes por especificação. "
    "Se este assert falhou, o padding mudou e os pacotes gerados "
    "serão incompatíveis com a outra ponta.");
static_assert(offsetof(CabecalhoDeRede, comprimento) == 2);
static_assert(std::is_trivially_copyable_v<CabecalhoDeRede>);
\end{lstlisting}

O subprojeto mostra o mesmo \texttt{struct} com e sem \texttt{\#pragma pack(1)},
e os tamanhos diferem --- que é precisamente o acidente que os três
\emph{asserts} existem para transformar em erro de \emph{build} em vez de pacote
corrompido em produção.

% -----------------------------------------------------------------------------
\subsubsection{NTTP: quando o valor entra no tipo}
% -----------------------------------------------------------------------------

Um \emph{template} pode receber tipos (\texttt{template <typename T>}) e também
\textbf{valores} (\texttt{template <int N>}). O segundo caso chama-se
\textbf{NTTP}, de \emph{non-type template parameter}, e o ponto é que o valor
passa a fazer parte da identidade do tipo: \texttt{BufferFixo<int,4>} e
\texttt{BufferFixo<int,16>} são tipos \textbf{diferentes}, não são atribuíveis
um ao outro, e \texttt{capacidade()} não lê memória nenhuma.

A aplicação que mais rende é verificação dimensional, porque transforma uma
classe inteira de defeitos de execução em erro de compilação:

\begin{lstlisting}[style=bt, language=C++,
    caption={Álgebra de matrizes verificada pelo compilador, em
             \texttt{src/compile\_time/demo\_nttp.hpp}.}]
template <std::size_t Linhas, std::size_t Colunas>
class Matriz {
    // Só compila se as dimensões internas casarem, e o resultado tem o
    // tipo certo, derivado das entradas.
    template <std::size_t OutrasColunas>
    constexpr Matriz<Linhas, OutrasColunas>
    operator*(const Matriz<Colunas, OutrasColunas>& o) const;
};

Matriz<2,3> * Matriz<3,4> → Matriz<2,4>     // compila
Matriz<2,3> * Matriz<2,2> → NÃO COMPILA     // dimensões incompatíveis
\end{lstlisting}

Sem NTTP, as dimensões seriam campos e o erro só apareceria em execução ---
provavelmente como leitura fora dos limites, que é \gls{ub} silencioso.

Até o C++17, NTTP só admitia inteiros, \texttt{enum}, ponteiros e referências. O
C++20 ampliou para tipos de classe literais, o que permitiu parametrizar por
texto. A condição é o tipo ser \textbf{\emph{structural}}: todos os membros de
dados públicos, sem \texttt{private} ou \texttt{protected}, sem
\texttt{virtual}, e os membros também \emph{structural}. É por isso que
\classe{std::string} \textbf{não} pode ser NTTP --- tem membros privados e aloca
--- e é por isso que se escreve um \texttt{StringFixa} à mão:

\begin{lstlisting}[style=bt, language=C++,
    caption={\emph{Strings} como parâmetro de \emph{template} (C++20).}]
template <std::size_t N>
struct StringFixa {
    char dados[N]{};
    consteval StringFixa(const char (&lit)[N]) { std::copy_n(lit, N, dados); }
    constexpr std::string_view view() const { return {dados, N - 1}; }
    constexpr bool operator==(const StringFixa&) const = default;  // obrigatório
};
template <std::size_t N> StringFixa(const char (&)[N]) -> StringFixa<N>;

template <StringFixa Nome, typename T> class Campo { /* ... */ };

Campo<"largura", int> largura{1920};
static_assert(decltype(largura)::nome() == "largura");
\end{lstlisting}

O \texttt{operator==} padronizado não é decoração: a comparação estrutural é
requisito para servir de NTTP, porque o compilador precisa decidir se duas
instanciações são a mesma. O \texttt{N - 1} corta o
\texttt{'\textbackslash 0'} do literal. E a consequência é que
\texttt{Campo<"largura",int>} e \texttt{Campo<"altura",int>} são tipos
\textbf{distintos}: trocá-los é erro de compilação, não erro de lógica. Vale
mencionar ainda o \texttt{template <auto Valor>}, que aceita qualquer constante
--- \texttt{int}, \texttt{double}, enumerador --- sem um \emph{template} por
tipo de valor.

\begin{armadilha}[cada valor distinto é uma instanciação nova]
NTTP tem um custo que não aparece em \emph{benchmark} nenhum de execução:
\texttt{BufferFixo<int,4>}, \texttt{<int,8>} e \texttt{<int,16>} são três
classes completas, com todos os métodos gerados três vezes. As consequências são
binário maior, compilação mais lenta e pressão sobre o cache de instruções num
laço que alterne entre instanciações. Use quando os valores forem
\textbf{poucos, fixos e conhecidos} e participarem da \textbf{correção} do
programa --- dimensões, unidades. Não use quando o valor vier de configuração,
arquivo ou entrada; e, se a intenção era apenas evitar uma alocação, o
instrumento certo é \texttt{std::span}. A técnica clássica para conter a
explosão é extrair para uma base não-\emph{template} tudo o que não depende do
NTTP.
\end{armadilha}

% -----------------------------------------------------------------------------
\subsubsection{Tabelas geradas na compilação}
% -----------------------------------------------------------------------------

Os recursos anteriores são mecanismo; este é o pagamento: substituir trabalho de
inicialização por dados que já nascem prontos dentro do binário. O exemplo
canônico é a tabela de 256 entradas do CRC-32, tradicionalmente gerada na
inicialização do programa ou colada à mão a partir de um arquivo gerado.

\begin{lstlisting}[style=bt, language=C++,
    caption={A tabela do CRC-32 calculada pelo compilador, em
             \texttt{src/compile\_time/demo\_tabelas.hpp}.}]
constexpr std::array<std::uint32_t, 256> gerar_tabela_crc32() {
    constexpr std::uint32_t polinomio = 0xEDB88320u;
    std::array<std::uint32_t, 256> tabela{};
    for (std::uint32_t i = 0; i < 256; ++i) {
        std::uint32_t c = i;
        for (int k = 0; k < 8; ++k)
            c = (c & 1u) ? (polinomio ^ (c >> 1)) : (c >> 1);
        tabela[i] = c;
    }
    return tabela;
}

// `inline constexpr` numa variável de header: uma única instância no
// programa, sem violar a ODR. Antes do C++17 isso exigiria um .cpp.
inline constexpr auto TABELA_CRC32 = gerar_tabela_crc32();

static_assert(TABELA_CRC32[255] == 0x2D02EF8Du);
\end{lstlisting}

A ODR --- \emph{One Definition Rule} --- exige que cada entidade tenha
exatamente uma definição no programa inteiro. Um \emph{header} incluído por dez
\texttt{.cpp} definiria a variável dez vezes; \texttt{inline} instrui o
\emph{linker} a fundir as definições numa só.

\paragraph{A tabela vai mesmo para o binário --- verificado.} Em vez de afirmar
que a tabela está no binário, o subprojeto manda inspecionar o binário:

\begin{lstlisting}[style=bt, language=bash,
    caption={Medição em \texttt{src/compile\_time/}: onde as tabelas foram parar
             no executável.}]
$ objdump -t dist/bin/compile_time | grep -i TABELA
.rodata      0x400  TABELA_CRC32
.data.rel.ro 0x0d8  TABELA_HTTP
\end{lstlisting}

O resultado contraria a expectativa ingênua de que as duas estariam no mesmo
lugar, e a razão é instrutiva. \texttt{TABELA\_CRC32} contém apenas
\texttt{uint32\_t} --- valores puros, gravados literalmente no arquivo e
mapeados somente-leitura: vão para \texttt{.rodata}, ocupando
$0x400 = 1024$ bytes, exatamente $256 \times 4$. \texttt{TABELA\_HTTP} contém
\classe{std::string\_view}, que são \textbf{ponteiros}: o endereço final só é
conhecido quando o \emph{loader} mapeia o binário, então ela precisa de correção
em tempo de carga e vai para \texttt{.data.rel.ro} --- somente-leitura
\emph{após} a relocação. O que mais importa é o que \emph{não} apareceu: nenhuma
das duas está em \texttt{.bss}, a seção de dados zerados que seria preenchida
por código de inicialização. Se estivessem lá, a promessa de ``calculado na
compilação'' seria falsa.

Desde o C++20 \texttt{std::sort} é \texttt{constexpr}, o que permite escrever a
tabela fora de ordem, deixar o compilador ordenar, e usar busca binária em
execução:

\begin{lstlisting}[style=bt, language=C++,
    caption={A pré-condição da busca binária, provada em vez de testada.}]
constexpr auto gerar_tabela_http() {
    std::array<EntradaHttp, 9> t{{ /* escrita fora de ordem */ }};
    std::sort(t.begin(), t.end());          // constexpr desde C++20
    return t;
}
inline constexpr auto TABELA_HTTP = gerar_tabela_http();

static_assert(std::is_sorted(TABELA_HTTP.begin(), TABELA_HTTP.end()));
static_assert(descrever_http(404) == "Not Found");
static_assert(descrever_http(999) == "Unknown");
\end{lstlisting}

Aplicar busca binária a uma sequência não ordenada é \gls{ub} silencioso
(\cref{sec:stl}). Aqui a pré-condição não é testada: é \emph{provada}. Essa
diferença é o argumento inteiro da seção.

\classe{std::map} não é \texttt{constexpr}, porque aloca nós. Para um mapeamento
pequeno e fixo, um arranjo ordenado com busca binária ganha nas três dimensões
(\cref{tab:mapa-fixo}) --- inclusive em velocidade, porque é contíguo, com a
mesma lógica de cache do \cref{sec:stl}.

\begin{table}[htbp]
\centering
\caption{Para poucas dezenas de entradas fixas, o arranjo ordenado ganha em tudo.}
\label{tab:mapa-fixo}
\footnotesize
\begin{tabular}{@{}lccc@{}}
\toprule
\textbf{Estrutura} & \texttt{constexpr} & \textbf{Aloca} & \textbf{Busca} \\
\midrule
\texttt{std::map}            & não          & sim          & $O(\log n)$ \\
\texttt{std::unordered\_map} & não          & sim          & $O(1)$ médio \\
\textbf{arranjo ordenado}    & \textbf{sim} & \textbf{não} & $O(\log n)$ \\
\bottomrule
\end{tabular}
\end{table}

\begin{armadilha}[uma API \texttt{constexpr} robusta usa índices, não ponteiros]
Esta lição só aparece escrevendo o código. A forma natural de devolver ``achou
ou não'' é um ponteiro, e depois \texttt{return p ? *p : padrao;}. Acontece que
converter um ponteiro-para-membro-de-objeto-\texttt{constexpr} em \texttt{bool}
\textbf{não é expressão constante} em todos os modos do GCC: o mesmo código
compila normalmente e falha sob \texttt{-fsanitize}, com a avaliação constante
recusada. O \texttt{MapaFixo} do subprojeto devolve, por isso, o
\textbf{índice} --- ou \texttt{N}, se não achou. Aritmética de índice é sempre
expressão constante, e por isso é a base mais robusta. É um caso em que a
medição --- aqui, um \emph{build} que falha --- contrariou o que o projeto
``natural'' recomendava.
\end{armadilha}

\begin{regra}[o custo é pago pelo compilador, e ele é lento]
O interpretador de \texttt{constexpr} do GCC é ordens de grandeza mais lento que
código nativo, e tem teto (\texttt{-fconstexpr-ops-limit}, 33,5 milhões de
passos). \textbf{Não} vale a pena quando a tabela é grande e o \emph{build} fica
lento, quando os dados mudam com frequência --- um arquivo lido em execução é
mais fácil de manter --- ou quando o cálculo já era trivial na inicialização.
\textbf{Vale muito} quando a tabela é pequena ou média e imutável (CRC, seno,
códigos de erro), quando se quer \emph{provar} uma propriedade em vez de
testá-la, em ambiente sem \emph{heap} (embarcado, \emph{kernel}), e para
eliminar inicialização estática dinâmica. A forma de decidir é medir:
\texttt{time g++ -std=c++23 -c src/compile\_time/main.cpp}.
\end{regra}

% =============================================================================
\subsection{O inventário medido do C++23}
\label{sec:cpp23-medido}
% =============================================================================

O subprojeto \projeto{cpp23} é diferente de todos os outros: além de
\emph{ensinar} os recursos do C++23, ele \textbf{mede} quais estão disponíveis
no compilador local e imprime o resultado. A tabela de suporte do
\texttt{src/ROADMAP.md} é gerada a partir dessa saída.

\begin{porque}[``o C++23'' não é algo que um compilador tem ou não tem]
Cada recurso do padrão chega numa versão diferente do compilador, e a biblioteca
costuma atrasar em relação à linguagem --- são equipes e cronogramas distintos.
Perguntar ``este compilador suporta C++23?'' é uma pergunta mal formulada: a
pergunta útil é ``suporta \emph{este recurso}?'', e ela tem resposta mecânica.
\end{porque}

Uma \textbf{\emph{feature-test macro}} é uma macro que o padrão obriga a
implementação a definir quando um recurso está disponível, com valor igual à
data de adoção no formato AAAAMM. Não são convenção de compilador --- são parte
do padrão, o que as torna a única forma portátil de escrever código que se
adapta. As de \textbf{linguagem} começam com \texttt{\_\_cpp\_} e são
predefinidas pelo compilador; as de \textbf{biblioteca} começam com
\texttt{\_\_cpp\_lib\_} e vivem em \cabecalho{version}, um \emph{header} que não
declara nada além delas. O valor numérico importa, porque recursos ganham
revisões: \texttt{\#if \_\_cpp\_lib\_optional >= 202110} testa especificamente
se as operações monádicas estão presentes.

\subsubsection{O inventário, e as seis entradas corrigidas}

\begin{lstlisting}[style=bt, language=bash,
    caption={Saída de \texttt{src/cpp23/}: o que este GCC 13.3 tem de
             \emph{linguagem}.}]
if consteval                    ✓ 202106
operator[](i, j, k)             ✓ 202211
static operator()               ✓ 202207
sufixo uz / z                   ✓ 202011
\N{NOME DO CARACTERE}           ✓ 202207
deducing this                   ✗ ausente     ← a mais sentida
\end{lstlisting}

Na biblioteca estão \textbf{presentes} \texttt{std::expected},
\texttt{std::byteswap}, \texttt{std::to\_underlying}, \texttt{std::unreachable},
\texttt{views::zip}, \texttt{string::contains} e \texttt{std::stacktrace} ---
este último com uma ressalva grave, tratada adiante. Estão \textbf{ausentes}
\texttt{std::print}/\texttt{println}, \texttt{std::mdspan},
\texttt{std::generator}, \texttt{ranges::to<C>()}, \texttt{std::flat\_map},
\texttt{views::enumerate}, \texttt{ranges::starts\_with} e
\texttt{std::out\_ptr}.

E aqui está o resultado que ensina método. A tabela de suporte do repositório
foi escrita, antes da medição, a partir de leitura de documentação. Quando o
programa passou a medir, \textbf{seis entradas estavam erradas}
(\cref{tab:presumido-medido}) --- todas na mesma direção: a de
\textbf{subestimar} o que já estava disponível. Ou seja, o repositório estava
reimplementando à mão coisas que a biblioteca já oferecia. O viés tem
explicação: documentação de terceiros e tabelas de compatibilidade envelhecem, e
a memória de quando um recurso ``ainda não existia'' persiste anos depois de
deixar de ser verdade. A macro não envelhece, porque é o próprio compilador
respondendo.

\begin{table}[htbp]
\centering
\caption{O que a documentação sugeria contra o que a medição encontrou no GCC 13.3.}
\label{tab:presumido-medido}
\footnotesize
\begin{tabular}{@{}lll@{}}
\toprule
\textbf{Recurso} & \textbf{Presumido} & \textbf{Medido} \\
\midrule
\texttt{std::stacktrace}               & ausente & \textbf{presente} (com \texttt{-lstdc++\_libbacktrace}) \\
\texttt{views::zip}                    & ausente & \textbf{presente} \\
\texttt{ranges::fold\_left}            & ausente & \textbf{presente} \\
\texttt{views::chunk} / \texttt{slide} & ausente & \textbf{presente} \\
\texttt{std::move\_only\_function}     & ausente & \textbf{presente} \\
\texttt{std::spanstream}               & ausente & \textbf{presente} \\
\bottomrule
\end{tabular}
\end{table}

\begin{regra}[não presuma: meça]
Reproduzir o inventário em qualquer \emph{toolchain} custa duas linhas:
\begin{lstlisting}[style=bt, language=bash, frame=none,
                   xleftmargin=0pt, xrightmargin=0pt]
echo '#include <version>' > /tmp/v.cpp
g++ -std=c++23 -dM -E /tmp/v.cpp | grep __cpp_lib | sort
\end{lstlisting}
Isso despeja todas as macros de biblioteca definidas, com os valores. É a
resposta autoritativa, e é local.
\end{regra}

\begin{armadilha}[a macro diz que a API existe, não que ela linka]
Esta é a lição mais útil do subprojeto inteiro. O
\texttt{\_\_cpp\_lib\_stacktrace} está definido neste GCC, e o \emph{header}
existe. Mesmo assim:
\begin{lstlisting}[style=bt, language=bash, frame=none,
                   xleftmargin=0pt, xrightmargin=0pt]
$ g++ -std=c++23 usa_stacktrace.cpp
undefined reference to `__glibcxx_backtrace_simple'

$ g++ -std=c++23 usa_stacktrace.cpp -lstdc++_libbacktrace
frames=5   ✓
\end{lstlisting}
A \emph{feature-test macro} responde sobre a \textbf{interface}, que é o que o
compilador enxerga; ela não sabe nada sobre bibliotecas de link. Sempre há um
passo de verificação além da macro, e ele é \textbf{compilar de verdade}. O
subprojeto opta por não usar \texttt{std::stacktrace} justamente porque
\texttt{libstdc++\_libbacktrace} não está em toda distribuição, e exigi-la
quebraria o \emph{build} de quem clonar o repositório.
\end{armadilha}

O padrão que permite usar um recurso novo sem quebrar em \emph{toolchains}
antigos é o \texttt{\#if} sobre a macro, com um caminho alternativo equivalente:

\begin{lstlisting}[style=bt, language=C++,
    caption={Três implementações escolhidas em compilação, em
             \texttt{src/cpp23/demo\_inventario.hpp}.}]
constexpr std::uint32_t trocar_bytes(std::uint32_t v) {
#if defined(__cpp_lib_byteswap) && __cpp_lib_byteswap >= 202110L
    return std::byteswap(v);
#elif defined(__GNUC__)
    return __builtin_bswap32(v);
#else
    return ((v & 0x000000FFu) << 24) | ((v & 0x0000FF00u) << 8) |
           ((v & 0x00FF0000u) >> 8)  | ((v & 0xFF000000u) >> 24);
#endif
}

static_assert(trocar_bytes(0x12345678u) == 0x78563412u);
\end{lstlisting}

Três regras governam esse padrão. \textbf{Um:} teste a \emph{macro}, não a
versão do compilador --- \texttt{\#if \_\_GNUC\_\_ >= 13} erra em Clang e em
MSVC. \textbf{Dois:} compare com o \emph{valor}, não apenas \texttt{\#ifdef},
porque recursos ganham revisões. \textbf{Três:} o caminho alternativo precisa
ser \emph{equivalente}, não aproximado --- e o \texttt{static\_assert} da última
linha é o que verifica isso.

\subsubsection{O que a linguagem ganhou}

Além do \texttt{if consteval} já tratado, quatro adições merecem registro.

O \textbf{subscript multidimensional} --- \texttt{operator[](i, j, k)} ---
finalmente existe, e funciona inclusive em \texttt{constexpr}. Mas vem com uma
armadilha de compatibilidade rara: uma sintaxe reciclada.

\begin{armadilha}[\texttt{a[i, j]} mudou de significado entre C++20 e C++23]
Em C++20, \texttt{a[i, j]} era o operador \textbf{vírgula} dentro dos colchetes:
avalia \texttt{i}, \textbf{descarta} o resultado, e indexa por \texttt{j}. O
C++20 depreciou esse uso e o C++23 reciclou a sintaxe. A consequência é que
código antigo com \texttt{a[i, j]} \textbf{muda de comportamento} ao ser
recompilado em C++23 --- silenciosamente, se o tipo de \texttt{a} tiver o
operador novo. Em arranjos crus a vírgula continua valendo, porque arranjos não
podem ter \texttt{operator[]} próprio. Compile com \texttt{-Wcomma-subscript}
para o GCC avisar.
\end{armadilha}

\texttt{static operator()} e \texttt{static operator[]} atendem ao caso do
\emph{functor} sem estado, que não precisa do \texttt{this}: marcar a chamada
como \texttt{static} remove o parâmetro implícito --- em código genérico muito
quente, é um argumento a menos por invocação. É por isso que os adaptadores de
\cabecalho{ranges} do C++23 usam essa forma: são chamados milhões de vezes e
nenhum tem estado. Note que \texttt{static} e \texttt{const} são mutuamente
exclusivos aqui: não há \texttt{this} para ser \texttt{const}.

Os \textbf{sufixos \texttt{uz} e \texttt{z}} resolvem a ausência de sufixo de
literal para \texttt{std::size\_t}, que obrigava a \texttt{static\_cast}
decorativos ou gerava avisos de comparação com sinais diferentes. Agora
\texttt{0uz} já é \texttt{size\_t} e \texttt{0z} é \texttt{std::ptrdiff\_t}:
\texttt{for (auto i = 0uz; i < v.size(); ++i)} não gera
\texttt{-Wsign-compare}.

\texttt{[[assume]]} é uma promessa ao otimizador: ``esta condição é verdadeira,
pode contar com isso''. Com ela, o compilador pode eliminar verificações,
escolher instruções mais estreitas e remover ramos.

\begin{armadilha}[\texttt{[[assume]]} é o oposto de \texttt{assert}]
\texttt{assert(x > 0)} \emph{verifica} e aborta se for falso;
\texttt{[[assume(x > 0)]]} \textbf{não verifica} e é \gls{ub} se for falso. E a
expressão \textbf{não é avaliada}: \texttt{[[assume(++i > 0)]]} não incrementa
\texttt{i}. Use com muita parcimônia, e sempre em par com um \texttt{assert} em
\emph{debug} --- assim o \emph{build} de depuração \textbf{prova} o que o de
\emph{release} \textbf{assume}.
\end{armadilha}

Por fim, os escapes nomeados:
\texttt{"\textbackslash N\{RIGHTWARDS ARROW\}"} usa o nome oficial do caractere
no Unicode em vez do número, e o compilador valida o nome --- um nome errado é
erro de compilação, não um caractere errado descoberto na tela.

\subsubsection{O que a biblioteca ganhou}

\texttt{std::to\_underlying} substitui
\texttt{static\_cast<std::underlying\_type\_t<Codigo>>(c)}. O ganho não é
digitar menos: é \textbf{não repetir o tipo}. Se \texttt{Codigo} mudar de
\texttt{uint16\_t} para \texttt{uint32\_t}, o \texttt{static\_cast} explícito
continua compilando e \textbf{trunca}; o \texttt{to\_underlying} acompanha
sozinho. Ele não faz o caminho de volta --- \texttt{int} para \texttt{enum
class} continua exigindo \texttt{static\_cast}, e continua sem validar que o
valor corresponde a um enumerador existente.

\texttt{std::unreachable} resolve um aviso legítimo: um \texttt{switch} que
cobre todos os enumeradores ainda faz o compilador dizer \emph{``control reaches
end of non-void function''}, porque um \texttt{enum class} \emph{pode} guardar
um valor fora da lista, obtido por \texttt{static\_cast}. Das três saídas, só
uma diz a verdade: \texttt{return \{\};} esconde o defeito atrás de um valor
falso; \texttt{throw} custa e trata um caso impossível;
\texttt{std::unreachable()} informa o compilador do que realmente se sabe. Ele é
equivalente a \texttt{[[assume(false)]]}, com a mesma consequência: se a
premissa estiver errada, é \gls{ub}.

\texttt{views::zip} preenche o buraco mais visível de \cabecalho{ranges} no
C++20 --- percorrer sequências em paralelo:

\begin{lstlisting}[style=bt, language=C++,
    caption={\texttt{zip} de três sequências, com \emph{structured bindings}.}]
for (const auto& [nome, idade, salario] : std::views::zip(nomes, idades, salarios))
    fmt::print("{:<7} {:>3} anos  R$ {:>8.2f}\n", nome, idade, salario);
\end{lstlisting}

Ele para no \textbf{menor} \emph{range}, o que elimina o \texttt{min} das formas
antigas com índice --- fácil de esquecer. Como \texttt{zip} referencia em vez de
copiar, valem as mesmas regras de tempo de vida das \emph{views}
(\cref{sec:ranges}). E o substituto para o \texttt{views::enumerate} ausente sai
de graça: \texttt{views::zip(views::iota(0uz), v)} funciona porque
\texttt{iota(0uz)} é infinito e o \texttt{zip} para no menor.

\classe{std::function} exige que o chamável seja \textbf{copiável}, mesmo quando
ninguém vai copiá-lo --- é uma exigência do tipo, não do uso. Isso impede
guardar um \emph{lambda} que capturou um \classe{std::unique\_ptr}, uma
\texttt{std::thread} ou um \texttt{std::promise}:

\begin{lstlisting}[style=bt, language=C++,
    caption={O chamável que possui um recurso.}]
std::function<std::string()> f = [r = std::move(ptr)] { return *r; };            // ✗
std::move_only_function<std::string()> f = [r = std::move(ptr)] { return *r; };  // ✓
\end{lstlisting}

Onde isso morde na prática é a fila de tarefas de um \emph{thread pool}: a
tarefa captura um \texttt{promise}, que é \emph{move-only}
(\cref{sec:concurrency}). E a API nova ainda qualifica a chamada ---
\texttt{void() const}, \texttt{void() \&\&}, \texttt{void() noexcept} --- coisas
que \texttt{std::function} nunca teve.

\texttt{std::spanstream} é o \emph{stream} sobre um \emph{buffer} que já existe.
Um \texttt{stringstream} sempre aloca; para analisar um \emph{buffer} pronto,
isso é uma cópia inteira desnecessária. O antigo \texttt{strstream} resolvia e
foi depreciado porque a API não deixava claro quem era dono do \emph{buffer}; o
\texttt{span} resolve isso por construção, porque \textbf{nunca} é dono.

\begin{armadilha}[o \texttt{- 1} do \texttt{span} sobre um \texttt{char[]} não é detalhe]
\begin{lstlisting}[style=bt, language=C++, frame=none,
                   xleftmargin=0pt, xrightmargin=0pt]
char dados[] = "42 3.14 texto";
std::ispanstream entrada{std::span<char>{dados, std::size(dados) - 1}};
\end{lstlisting}
\texttt{std::span\{dados\}} sobre um \texttt{char[]} inclui o
\texttt{'\textbackslash 0'} final --- o arranjo tem 14 bytes, não 13. Sem o
\texttt{- 1}, o terminador vira um caractere de dado e a última \emph{string}
lida sai com um byte a mais. O sintoma é um tamanho errado que \emph{parece}
certo quando impresso.
\end{armadilha}

Por fim, duas \emph{views} de janela que parecem iguais e não são:

\begin{lstlisting}[style=bt, language=bash,
    caption={Saída de \texttt{src/cpp23/}: janelas disjuntas contra deslizantes.}]
chunk(3): [1 2 3] [4 5 6] [7 8 9] [10]        ← disjuntos, o último pode ser menor
slide(3): [1 2 3] [2 3 4] [3 4 5] ...          ← deslizante, todos do mesmo tamanho
\end{lstlisting}

\texttt{chunk} avança de $N$ em $N$ e serve para processar em lotes;
\texttt{slide} avança de 1 em 1 e serve para média móvel --- que fica em uma
linha, combinando \texttt{slide(3)} com \texttt{ranges::fold\_left}. Antes do
C++23, os dois eram laços com índice e aritmética de fronteira fácil de errar no
último bloco.

\subsubsection{O que falta, e o substituto de cada coisa}

Toda ausência vem com uma alternativa que funciona hoje
(\cref{tab:cpp23-faltando}). O objetivo não é lamentar: é escrever código que já
resolve o problema e que migra sem dor quando o \emph{toolchain} acompanhar. O
caso do \texttt{std::print} merece nota: ele \textbf{nasceu} da biblioteca
\texttt{fmt} --- a proposta P2093 foi escrita pelo autor dela, e a API é
idêntica de propósito. É por isso que este repositório usa \texttt{fmt} em todos
os subprojetos: quando o \texttt{std::print} chegar, a migração é trocar o
prefixo.

\begin{table}[htbp]
\centering
\caption{O que o GCC 13.3 não tem, e o que usar no lugar.}
\label{tab:cpp23-faltando}
\footnotesize
\begin{tabular}{@{}ll@{}}
\toprule
\textbf{Ausente} & \textbf{Substituto que funciona hoje} \\
\midrule
\texttt{std::print} / \texttt{println} & \texttt{fmt::print} / \texttt{fmt::println} --- API idêntica \\
\texttt{std::mdspan}          & \texttt{operator[](i,j)} próprio; ou \texttt{kokkos/mdspan} \\
\texttt{ranges::to<C>()}      & \texttt{std::vector v(view.begin(), view.end())} \\
\texttt{views::enumerate}     & \texttt{views::zip(views::iota(0uz), v)} \\
\texttt{std::generator}       & \texttt{promise\_type} à mão (\cref{sec:coroutines}) \\
\emph{deducing this}          & \emph{y-combinator} (\cref{sec:lambdas}); CRTP \\
\texttt{std::flat\_map}       & arranjo ordenado (\cref{sec:ct-fundo}) \\
\texttt{ranges::starts\_with} & \texttt{std::equal} \\
\bottomrule
\end{tabular}
\end{table}

\begin{armadilha}[o substituto de \texttt{ranges::to} exige um \emph{common range}]
Construir um \texttt{vector} a partir de \texttt{view.begin()} e
\texttt{view.end()} só funciona se os dois forem do \emph{mesmo tipo} --- um
\emph{common range}. \emph{Views} que terminam em sentinela (\cref{sec:stl})
precisam de \texttt{views::common} antes, ou de um auxiliar genérico de seis
linhas que percorre o \emph{range} e empurra num \texttt{vector}.
\end{armadilha}

\begin{regra}[a régua para adotar um recurso novo]
A pergunta não é ``o C++23 tem?'', é ``o meu \emph{toolchain} \textbf{mínimo}
tem?''. Na ordem: (1)~qual é o compilador \textbf{mais antigo} que precisa
compilar este código --- a integração contínua, a distribuição LTS, o cliente?
(2)~a \emph{feature-test macro} está definida nele? (3)~ela \textbf{linka} sem
\emph{flags} extras? Se as três respostas forem sim, use direto. Se a segunda ou
a terceira falhar, o padrão é o \texttt{\#if} com caminho alternativo
equivalente.
\end{regra}

% =============================================================================
\subsection{Módulos nomeados}
\label{sec:modulos}
% =============================================================================

O último subprojeto enfrenta a maior mudança estrutural que o C++20 propôs: a
substituição do \texttt{\#include}.

\begin{porque}[o \texttt{\#include} é uma cópia de texto, e todos os problemas vêm daí]
O pré-processador não entende C++. Ele lê \texttt{\#include <vector>} e
substitui a linha pelo \emph{conteúdo textual} do arquivo, recursivamente. Cinco
consequências decorrem disso. \textbf{Compilação lenta:} o mesmo \emph{header} é
reprocessado do zero em \emph{cada} \texttt{.cpp} --- um projeto com 200 arquivos
que incluem \cabecalho{string} analisa \cabecalho{string} 200 vezes.
\textbf{Macros vazando:} um \texttt{\#define} feito antes do \texttt{\#include}
afeta o texto incluído, e vice-versa, sem barreira. \textbf{A ordem passa a
importar:} o \emph{header} B só compila se A vier antes, e isso não está escrito
em lugar nenhum. \textbf{\emph{Include guards}:} necessários porque a inclusão é
textual. \textbf{Tudo é público:} quem inclui A recebe também tudo que A
incluiu. Módulos resolvem os cinco de uma vez --- ao custo de uma \textbf{ordem
de compilação obrigatória}, que é o que torna a adoção difícil.
\end{porque}

Uma \textbf{interface de módulo} é um arquivo cujo conteúdo significativo começa
com \texttt{export module NOME;}. A extensão \texttt{.cppm} é convenção --- o
Clang a reconhece, o GCC não.

\begin{lstlisting}[style=bt, language=C++,
    caption={A estrutura completa, em \texttt{src/modules/geometria.cppm}.}]
module;                        // ← início do "global module fragment"
#include <cmath>               //   o ÚNICO lugar para #include numa interface

export module geometria;       // ← a declaração que define o módulo

namespace detalhe {            // não exportado → invisível para quem importa
    double quadrado(double x) { return x * x; }
}

export namespace geometria {   // exportado → a interface pública
    struct Ponto { double x{}, y{}; };

    double distancia(const Ponto& a, const Ponto& b) {
        return std::sqrt(detalhe::quadrado(b.x - a.x) +
                         detalhe::quadrado(b.y - a.y));
    }

    template <typename T>      // templates funcionam: o .gcm guarda a árvore
    T maior(T a, T b) { return a > b ? a : b; }

    inline constexpr double PI = 3.14159265358979323846;
}
\end{lstlisting}

O \textbf{\emph{global module fragment}} é o trecho entre \texttt{module;} e
\texttt{export module}: existe porque \emph{headers} antigos ainda precisam ser
incluídos, e é o único lugar de uma interface onde um \texttt{\#include} pode
aparecer. A diferença de fundo em relação a um \emph{header} é que o módulo é
compilado \textbf{uma vez}, e o consumidor lê a forma binária do resultado ---
no GCC, um arquivo \texttt{.gcm}. Templates funcionam normalmente, ao contrário
do que a intuição sugere: o \texttt{.gcm} guarda a árvore sintática, não texto.

\begin{table}[htbp]
\centering
\caption{O que muda entre \emph{header} e módulo.}
\label{tab:header-modulo}
\footnotesize
\begin{tabular}{@{}lll@{}}
\toprule
 & \textbf{Header} & \textbf{Módulo} \\
\midrule
compilação            & reprocessado em cada TU         & compilado \textbf{uma} vez \\
ordem                 & irrelevante                     & \textbf{obrigatória} \\
isolamento            & tudo vaza                       & só o que tem \texttt{export} \\
macros                & atravessam nos dois sentidos    & \textbf{não} atravessam \\
\emph{include guards} & necessários                     & desnecessários \\
templates             & precisam estar no \emph{header} & ficam no \texttt{.gcm} \\
\bottomrule
\end{tabular}
\end{table}

O demo prova o isolamento com um erro provocado de propósito.
\texttt{geometria.cppm} inclui \cabecalho{cmath} no \emph{global module
fragment} e usa \texttt{std::sqrt} internamente; o \texttt{main.cpp}
\textbf{não} inclui \cabecalho{cmath}:

\begin{lstlisting}[style=bt, language=bash,
    caption={Medição em \texttt{src/modules/}: o \texttt{\#include} do módulo
             não atravessa o \texttt{import}.}]
// em main.cpp:  double r = std::sqrt(16.0);
error: 'sqrt' is not a member of 'std'
\end{lstlisting}

Com um \emph{header}, \texttt{std::sqrt} estaria disponível ali sem ninguém ter
pedido --- e a poluição é \textbf{cumulativa}: um \emph{header} que inclui outro
que inclui outro acaba trazendo milhares de nomes. O mesmo vale para
\texttt{detalhe::quadrado}, definida sem \texttt{export}: é invisível no
consumidor, e o nome não colide com nada dele.

Um módulo pode ainda importar outro e repassá-lo:
\texttt{export import geometria;} em \texttt{formatacao.cppm} faz com que
\texttt{main.cpp}, que só escreve \texttt{import formatacao;}, enxergue
\texttt{geometria} também. E a ordem passa a importar: o \texttt{geometria.gcm}
precisa existir antes de \texttt{formatacao.cppm} compilar.

\subsubsection{Dois defeitos do GCC 13, encontrados escrevendo o demo}

Estas duas armadilhas custaram tempo real e não aparecem nos tutoriais. Ambas
foram verificadas no GCC 13.3 com \texttt{-fmodules-ts}.

\begin{armadilha}[\texttt{export} de variável global some com o \emph{global module fragment}]
Reduzido ao mínimo:
\begin{lstlisting}[style=bt, language=C++, frame=none,
                   xleftmargin=0pt, xrightmargin=0pt]
module;
#include <cmath>
export module m;
export constexpr double PI = 3.14;   // ← invisível no consumidor

// no consumidor:
import m;
int main() { return (int)PI; }       // error: 'PI' was not declared
\end{lstlisting}
\textbf{Sem} o \emph{global module fragment} --- removendo o \texttt{module;} e
o \texttt{\#include} do topo --- a mesma declaração funciona. Dentro de um
\texttt{namespace} exportado, funciona nos dois casos. É por isso que o
\texttt{PI} do subprojeto está dentro de \texttt{export namespace geometria}:
não é estilo, é necessidade.
\end{armadilha}

\begin{armadilha}[\texttt{constexpr} em escopo de \emph{namespace} tem \emph{linkage} interna]
\begin{lstlisting}[style=bt, language=C++, frame=none,
                   xleftmargin=0pt, xrightmargin=0pt]
export namespace geometria {
    constexpr double PI = 3.14;          // ✗ 'PI' is not a member of 'geometria'
    inline constexpr double PI = 3.14;   // ✓
}
\end{lstlisting}
\emph{Linkage} é a propriedade que determina se um nome pode ser referenciado de
outra unidade de tradução. Uma variável \texttt{constexpr} em escopo de
\emph{namespace} tem \emph{linkage} \textbf{interna} por padrão, como um
\texttt{const}: cada TU teria a sua cópia, e o módulo não a exporta de verdade.
A regra é a mesma que vale para \emph{headers} desde o C++17; o que muda é o
sintoma. Num \emph{header}, a consequência é uma violação silenciosa da ODR ---
cada TU com o seu endereço; num módulo, é erro de compilação no consumidor. O
erro é preferível, mas só se souber lê-lo.
\end{armadilha}

\subsubsection{As armadilhas do \emph{toolchain}}

\begin{armadilha}[\texttt{.cppm} precisa de \texttt{-x c++} no GCC]
\begin{lstlisting}[style=bt, language=bash, frame=none,
                   xleftmargin=0pt, xrightmargin=0pt]
g++ -fmodules-ts -c geometria.cppm          # ✗ "linker input file unused"
g++ -fmodules-ts -x c++ -c geometria.cppm   # ✓
\end{lstlisting}
O GCC não reconhece a extensão \texttt{.cppm} e trata o arquivo como entrada de
\emph{linker}. A mensagem de erro não sugere a causa. O Clang reconhece.
\end{armadilha}

Compilar uma interface produz \textbf{dois} artefatos: o \texttt{.o}, com o
código, e o \texttt{.gcm}, com a interface, num diretório \texttt{gcm.cache/}
criado automaticamente. O \texttt{.gcm} \textbf{não entra no link} --- é
consumido em tempo de compilação pelos consumidores. E o tamanho dele é
revelador:

\begin{lstlisting}[style=bt, language=bash,
    caption={Medição em \texttt{src/modules/}: os artefatos binários de módulo.}]
$ make mostrar-cache
-rw-r--r-- 3.3K  gcm.cache/geometria.gcm
-rw-r--r-- 2.1M  gcm.cache/formatacao.gcm      ← inclui <string> transitivamente
\end{lstlisting}

A diferença de três ordens de grandeza tem explicação direta:
\texttt{formatacao} incluiu \cabecalho{string} no seu \emph{global module
fragment}, e o \texttt{.gcm} guarda a representação compilada de tudo aquilo. É
o mesmo custo de sempre --- só que agora pago uma vez e visível num arquivo, em
vez de pago em toda TU e invisível.

Além dos módulos nomeados, o GCC 13 suporta \emph{header units}: pré-compilar um
\emph{header} com
\texttt{g++ -fmodules-ts -x c++-system-header vector} e depois escrever
\texttt{import <vector>;} --- em linha própria, porque
\texttt{import <vector>; int main()} na mesma linha não compila. O
\texttt{import std;}, que traria a biblioteca padrão inteira como módulo,
\textbf{não} está disponível no GCC 13.

\subsubsection{Por que este subprojeto fica fora do Meson}

Esta é a única exceção do repositório: \projeto{modules} tem um
\texttt{Makefile} próprio e não participa do \texttt{make build}. A razão é
conteúdo didático, não conveniência.

\begin{porque}[a ordem de compilação é o problema não resolvido de módulos]
Um sistema de \emph{build} precisa descobrir a ordem de compilação
\textbf{antes} de compilar. Com \texttt{\#include} isso é fácil: o
\texttt{-MMD} do GCC emite a lista de arquivos incluídos e o Make ou o Ninja
constroem o grafo. Com \texttt{import}, não: a dependência é sobre o
\texttt{.gcm}, que \emph{só existe depois de compilar} o módulo que o produz ---
uma circularidade entre o planejamento e a execução do \emph{build}.

A saída é um \emph{scanner} que lê os fontes e extrai as declarações
\texttt{import} antes de qualquer compilação: o \texttt{-fdeps-format}
padronizado, que chegou no GCC 14, ou o \texttt{clang-scan-deps}. O suporte do
Meson a módulos com GCC 13 é experimental e frágil justamente porque o formato
que o GCC 13 emite não é o que o Meson espera. \textbf{É exatamente esse
problema de \emph{tooling} que atrasou a adoção de módulos por cinco anos}, e
não qualquer defeito da ideia.
\end{porque}

O \texttt{Makefile} de 40 linhas do subprojeto demonstra melhor que a
integração, porque a ordem fica explícita e escrita à mão:

\begin{lstlisting}[style=bt, language=bash,
    caption={As dependências que o Make não descobre sozinho, em
             \texttt{src/modules/Makefile}.}]
# regra para interfaces de módulo: o -x c++ é obrigatório
%.o: %.cppm
	$(CXX) $(CXXFLAGS) -x c++ -c $< -o $@

# um #include o Make rastreia com -MMD; um import NÃO
formatacao.o: geometria.o
main.o:       formatacao.o geometria.o
\end{lstlisting}

\begin{table}[htbp]
\centering
\caption{Estado da adoção de módulos, por \emph{toolchain}.}
\label{tab:adocao-modulos}
\footnotesize
\begin{tabular}{@{}lccc@{}}
\toprule
\textbf{Compilador} & \textbf{Módulos nomeados} & \textbf{\emph{Header units}} & \texttt{import std;} \\
\midrule
GCC 13      & sim (\texttt{-fmodules-ts}) & sim & não \\
GCC 14+     & sim                         & sim & não \\
Clang 16+   & sim                         & sim & parcial \\
MSVC 19.3x  & sim                         & sim & sim \\
\bottomrule
\end{tabular}
\end{table}

\begin{regra}[quando adotar módulos]
Só se o compilador mínimo \textbf{e} o sistema de \emph{build} suportarem --- e
o segundo é o gargalo real, não o primeiro. Hoje isso significa CMake 3.28 ou
superior com Ninja 1.11 ou superior. Para migrar um projeto grande, comece pelas
\textbf{folhas} do grafo de dependências: os módulos que não importam ninguém
são os únicos que podem ser convertidos sem tocar em mais nada. Constantes
exportadas: sempre \texttt{inline constexpr}, dentro de um \emph{namespace}
exportado. Para esconder um detalhe: simplesmente não o exporte.
\end{regra}

% =============================================================================
\subsection{Síntese}
\label{sec:vocab-ct-sintese}
% =============================================================================

As quatro seções defendem a mesma tese por caminhos diferentes:
\textbf{informação que está no tipo é verificada; informação que está numa
convenção, não}.

Os tipos de vocabulário movem para o tipo aquilo que antes era combinado por
comentário --- ``pode ser nulo'', ``o segundo parâmetro é o tamanho do
primeiro'', ``só leia o membro que a etiqueta indicar''. O ganho é que o
compilador passa a poder recusar o uso errado; o custo é aprender cinco tipos e
as armadilhas de tempo de vida de dois deles.

A computação em tempo de compilação move para o compilador aquilo que antes era
feito na inicialização --- tabelas, validação de literais, ordenação. O ganho
maior não é o tempo economizado na execução: é que o avaliador de expressões
constantes recusa \gls{ub}, transformando uma classe inteira de defeitos em erro
de \emph{build}. O custo é tempo de compilação, que é real e mensurável.

O subprojeto de C++23 acrescenta uma lição de método que vale além do assunto:
\textbf{meça em vez de supor}. Seis entradas de uma tabela escrita a partir de
documentação estavam erradas, todas subestimando o que já existia. E a
verificação não termina na \emph{feature-test macro}: o \texttt{std::stacktrace}
está declarado, tem \emph{header}, e não linka.

Os módulos são o único dos quatro assuntos em que a recomendação prática é
esperar. A ideia resolve os cinco problemas do \texttt{\#include} de uma vez, e
a implementação nos compiladores existe; o que ainda não existe, de forma
madura, é o suporte dos sistemas de \emph{build} ao problema de descobrir a
ordem de compilação antes de compilar. Os dois defeitos do GCC 13 documentados
aqui --- o \texttt{export} que some com o \emph{global module fragment} e o
\texttt{constexpr} sem \texttt{inline} --- são sintomas dessa juventude.
